% basic nastaveni
\documentclass[12pt]{article}
\setlength{\parindent}{0pt}
\setlength{\parskip}{1em}
\usepackage[utf8]{inputenc}
\usepackage[czech]{babel}
\usepackage[T1]{fontenc}
\usepackage{graphics}
\usepackage{caption}
\usepackage{hyperref}

% matematicke package
\usepackage{amsmath}
\usepackage{amsfonts}
\usepackage{amssymb}
\usepackage{geometry}
\geometry{margin=2.5cm}
\usepackage{pgfplots}
\pgfplotsset{compat=1.18}
\setcounter{secnumdepth}{0}
\title{Maturitní otázka číslo 1: Logická výstavba matematiky}
\author{}
\date{}

\begin{document}
\maketitle
\subsection{Základní pojmy}
\textbf{Axiom} je základní tvrzení, jehož pravdivost se nedokazuje a stavíme na ní dále daný důkaz
\textbf{Věta} je tvrzení, jehož pravdivost je potřeba ověřit
\textbf{Výrok} je věta u které má smysl ověřit, jestli je pravdivá. Výroky jsou oznamovací věty, nebo matematická tvrzení, například:

- Brno neexistuje

- $4 + 4 = 5$

\textbf{Výroková forma} je výraz, který obsahuje jednu či více proměnných a není výrokem, dokud za ně nedosadíme. Výroková forma je například $4 + x = 8$
\subsection{Logické spojky}
\textbf{Konjunkce} se značí $\wedge$ a v češtině znamená a/a zároveň.\\
\textbf{Disjunkce} se značí $\vee$ a znamená nebo\\
\textbf{Implikace} se značí $A \rightarrow B$, výrok A se nayývá předpoklad a výrok B se nazývá závěr (jestliže A, pak B)\\
\textbf{Ekvivalence} se značí $A \leftrightarrow B$, znamená A je ekvivalentní s B\\



\end{document}
